%% Angles associés : les trois symétries du cercle trigonométrique (leçon p. 99).
%% M est associé à x (ici x = pi/5, soit 36°) ; M1 = symétrique par rapport à l'axe des
%% abscisses (-x), M2 = symétrique par rapport à l'axe des ordonnées (pi - x),
%% M3 = symétrique par rapport à O (pi + x).
\documentclass[tikz,border=4pt]{standalone}
\usepackage{liaisons}
\begin{document}
\begin{tikzpicture}[x=2cm,y=2cm,
  axe/.style={-{Stealth[length=5pt]}, line width=0.6pt},
  cercle/.style={solideE, line width=1.1pt},
  rappel/.style={line width=0.5pt, dash pattern=on 3pt off 2pt, black!55},
  etiq/.style={font=\small, fill=white, inner sep=1.5pt}]

\def\cosx{0.809}
\def\sinx{0.5878}

%% ---------------- repère et cercle -------------------------------------
\draw[axe] (-1.3,0) -- (1.3,0);
\draw[axe] (0,-1.28) -- (0,1.28);
\draw[cercle] (0,0) circle (1);
\node[font=\small, anchor=north east, inner sep=2pt] at (-0.03,-0.03) {$O$};

%% ---------------- les trois segments de symétrie -----------------------
\draw[rappel] (\cosx,\sinx) -- (\cosx,-\sinx);      % symétrie / axe des abscisses
\draw[rappel] (\cosx,\sinx) -- (-\cosx,\sinx);      % symétrie / axe des ordonnées
\draw[rappel] (\cosx,\sinx) -- (-\cosx,-\sinx);     % symétrie / centre O
\draw[rappel] (-\cosx,\sinx) -- (-\cosx,-\sinx);
\draw[rappel] (\cosx,-\sinx) -- (-\cosx,-\sinx);

%% ---------------- arc de l'angle x -------------------------------------
\draw[solideA, line width=1pt] (0.34,0) arc (0:36:0.34);
\node[font=\small, text=solideA, anchor=west, inner sep=1.5pt] at (17:0.38) {$x$};

%% ---------------- les quatre points ------------------------------------
\draw[solideA, line width=1.2pt] (0,0) -- (\cosx,\sinx);
\fill[solideA] (\cosx,\sinx) circle (1.8pt);
\node[font=\small, anchor=west, inner sep=3pt] at (1.02*\cosx,1.08*\sinx) {$M\,(x)$};
\fill (\cosx,-\sinx) circle (1.7pt);
\node[font=\small, anchor=west, inner sep=3pt] at (1.02*\cosx,-1.08*\sinx) {$M_1\,(-x)$};
\fill (-\cosx,\sinx) circle (1.7pt);
\node[font=\small, anchor=east, inner sep=3pt] at (-1.02*\cosx,1.08*\sinx) {$M_2\,(\pi-x)$};
\fill (-\cosx,-\sinx) circle (1.7pt);
\node[font=\small, anchor=east, inner sep=3pt] at (-1.02*\cosx,-1.08*\sinx) {$M_3\,(\pi+x)$};

%% ---------------- abscisses et ordonnées lues sur les axes -------------
\node[etiq, anchor=north] at (\cosx,-0.02) {$\cos(x)$};
\node[etiq, anchor=north] at (-\cosx,-0.02) {$-\cos(x)$};
\node[etiq, anchor=east] at (-0.03,\sinx) {$\sin(x)$};
\node[etiq, anchor=east] at (-0.03,-\sinx) {$-\sin(x)$};
\end{tikzpicture}
\end{document}
